\documentclass[10pt,a4paper]{amsart}
\usepackage{amsmath,amssymb,amsthm,graphicx}
\usepackage{verbatim}
\newenvironment{reference}{\comment}{\endcomment}
%\newenvironment{reference}{}{}
\newenvironment{slogan}{\comment}{\endcomment}
\newenvironment{history}{\comment}{\endcomment}

% For commutative diagrams we use Xy-pic
\usepackage[all]{xy}

% We use 2cell for 2-commutative diagrams.
\xyoption{2cell}
% Theorem environments.
%
\theoremstyle{plain}
\newtheorem{theorem}[subsection]{Theorem}
\newtheorem{proposition}[subsection]{Proposition}
\newtheorem{lemma}[subsection]{Lemma}

\theoremstyle{definition}
\newtheorem{definition}[subsection]{Definition}
\newtheorem{example}[subsection]{Example}
\newtheorem{exercise}[subsection]{Exercise}
\newtheorem{situation}[subsection]{Situation}

\theoremstyle{remark}
\newtheorem{remark}[subsection]{Remark}
\newtheorem{remarks}[subsection]{Remarks}

\numberwithin{equation}{subsection}

% Macros
%
\def\lim{\mathop{\mathrm{lim}}\nolimits}
\def\colim{\mathop{\mathrm{colim}}\nolimits}
\def\Spec{\mathop{\mathrm{Spec}}}
\def\Hom{\mathop{\mathrm{Hom}}\nolimits}
\def\Ext{\mathop{\mathrm{Ext}}\nolimits}
\def\SheafHom{\mathop{\mathcal{H}\!\mathit{om}}\nolimits}
\def\SheafExt{\mathop{\mathcal{E}\!\mathit{xt}}\nolimits}
\def\Sch{\mathit{Sch}}
\def\Mor{\mathop{\mathrm{Mor}}\nolimits}
\def\Ob{\mathop{\mathrm{Ob}}\nolimits}
\def\Sh{\mathop{\mathit{Sh}}\nolimits}
\def\NL{\mathop{N\!L}\nolimits}
\def\CH{\mathop{\mathrm{CH}}\nolimits}
\def\proetale{{pro\text{-}\acute{e}tale}}
\def\etale{{\acute{e}tale}}
\def\QCoh{\mathit{QCoh}}
\def\Ker{\mathop{\mathrm{Ker}}}
\def\Im{\mathop{\mathrm{Im}}}
\def\Coker{\mathop{\mathrm{Coker}}}
\def\Coim{\mathop{\mathrm{Coim}}}

% Boxtimes
%
\DeclareMathSymbol{\boxtimes}{\mathbin}{AMSa}{"02}

%
% Macros for moduli stacks/spaces
%
\def\QCohstack{\mathcal{QC}\!\mathit{oh}}
\def\Cohstack{\mathcal{C}\!\mathit{oh}}
\def\Spacesstack{\mathcal{S}\!\mathit{paces}}
\def\Quotfunctor{\mathrm{Quot}}
\def\Hilbfunctor{\mathrm{Hilb}}
\def\Curvesstack{\mathcal{C}\!\mathit{urves}}
\def\Polarizedstack{\mathcal{P}\!\mathit{olarized}}
\def\Complexesstack{\mathcal{C}\!\mathit{omplexes}}
% \Pic is the operator that assigns to X its picard group, usage \Pic(X)
% \Picardstack_{X/B} denotes the Picard stack of X over B
% \Picardfunctor_{X/B} denotes the Picard functor of X over B
\def\Pic{\mathop{\mathrm{Pic}}\nolimits}
\def\Picardstack{\mathcal{P}\!\mathit{ic}}
\def\Picardfunctor{\mathrm{Pic}}
\def\Deformationcategory{\mathcal{D}\!\mathit{ef}}
\title{Minimal objects of the category of morphisms to a fixed object under a base change}
\author{The Stacks Project Authors}
\date{stacks-project:formal-defos.tex:local-snapshot-20261012; GNU FDL 1.2 or later}
\begin{document}
\maketitle
\noindent\textit{Editorial note.} All delivered mathematical text is the author's own, copied from the pinned Stacks Project chapter file formal-defos.tex; no statement or proof step was rewritten, reordered or removed.\par\medskip
\noindent\textit{Editorial note.} The counted proof is the author's complete argument: the reduction of a competing morphism in the target category to the source category through (S1) and the resulting identicalness of the image of the minimal object, and the lifting along a surjection with the control of the resulting object by minimality.\par\medskip
\noindent\textit{Editorial note.} The label lemma-smallest-where-descends names the counted statement, and the proof environment immediately following it in the source is delivered verbatim as its proof.\par\medskip
\noindent\textit{Editorial note.} The counted statement is carried by the author's own lemma environment in the chapter source.\par\medskip
\noindent\textit{Editorial note.} Pre-open gate: python operations/agent-authoring/tools/precard\_gate.py --source formal-defos.tex --label lemma-smallest-where-descends --start 4321 --end 4372 --local-id 28021 returned RESULT=OK (L28021 free, occupancy FREE, public interval clear). Candidate taken from sources/stacks-formal-defos/candidates-halfA.json.\par\medskip
\noindent\textit{Editorial note.} The statement and proof are copied from the pinned chapter file; the cited inputs are kept as external pointers with their canonical Stacks tags.\par\medskip
\noindent\textit{Allowed standard prerequisites (editorial).}\begin{itemize}
\item Construction of the pushforward of a morphism along a ring map. Same chapter, the construction of f\_* of a morphism y -> x along a ring map f, using the (S1) property to produce the displayed square. Supplies the construction of f\_* and the displayed diagram lying over B' -> f(B') inside B -> A.
\item (S1) property of the category. Same chapter, the statement of (S1) as the existence of objects over a fibre product of Artinian quotients. Supplies the objects y'' and y' used in the two parts of the proof.
\item Existence of minimal objects dominated by a given one. Same chapter, the remark that every object of the category of morphisms admits a morphism from a minimal one, by boundedness of the lengths. Supplies the minimal object y'' -> y used at the end of the second part.
\end{itemize}\medskip
\begin{lemma}
\label{lemma-smallest-where-descends}
Let $\mathcal{F}$ be a category cofibred in groupoids over
$\mathcal{C}_\Lambda$ which has (S1).
\begin{enumerate}
\item For $y \to x$ in $\mathcal{F}$ a minimal
object in $\mathcal{S}_y$ maps to a minimal object of $\mathcal{S}_x$.
\item For $y \to x$ in $\mathcal{F}$ lying over a surjection
$f : B \to A$ in $\mathcal{C}_\Lambda$ every minimal object
of $\mathcal{S}_x$ is the image of a minimal object of
$\mathcal{S}_y$.
\end{enumerate}
\end{lemma}

\begin{proof}
Proof of (1). Say $y \to x$ lies over $f : B \to A$. Let $y' \to y$
lying over $B' \subset B$ be a minimal object of $\mathcal{S}_y$. Let
$$
\vcenter{
\xymatrix{
y' \ar[d] \ar[r] & x' \ar[d] \\
y \ar[r] & x
}
}
\quad\text{lying over}\quad
\vcenter{
\xymatrix{
B' \ar[d] \ar[r] & f(B') \ar[d] \\
B \ar[r] & A
}
}
$$
be as in the construction of $f_*$ above. Suppose that
$(x'' \to x) \to (x' \to x)$ is a morphism of $\mathcal{S}_x$
with $x'' \to x'$ lying over $A'' \subset f(B')$. By (S1)
there exists $y'' \to y'$ lying over $B' \times_{f(B')} A'' \to B'$.
Since $y' \to y$ is minimal we conclude that
$B' \times_{f(B')} A'' \to B'$ is an isomorphism, which implies that
$A'' = f(B')$, i.e., $x' \to x$ is minimal.

\medskip\noindent
Proof of (2). Suppose $f : B \to A$ is surjective and $y \to x$ lies over $f$.
Let $x' \to x$ be a minimal object of $\mathcal{S}_x$ lying over $A' \subset A$.
By (S1) there exists $y' \to y$ lying over
$B' = f^{-1}(A') = B \times_A A' \to B$ whose image in $\mathcal{S}_x$ is
$x' \to x$. So $f_*(y' \to y) = x' \to x$.
Choose a morphism $(y'' \to y) \to (y' \to y)$ in
$\mathcal{S}_y$ with $y'' \to y$ a minimal object (this is possible by
the remark on lengths above the lemma). Then $f_*(y'' \to y)$ is an
object of $\mathcal{S}_x$ which maps to $x' \to x$ (by functoriality of
$f_*$) hence is isomorphic to $x' \to x$ by minimality of $x' \to x$.
\end{proof}

\end{document}
